\Askhsh{33998_SOLUTION}{1}
\begin{ylh}{1}
    \item [Ύλη από εικόνα]
\end{ylh}
ΛΥΣΗ
\begin{alist}
    \item [α)] Έχουμε:
    \[
        \int_0^2 5 \cdot \left(\frac{4}{5}\right)^t \cdot \ln\frac{4}{5} dt = 5 \int_0^2 \left(\frac{4}{5}\right)^t \cdot \ln\frac{4}{5} dt = 5 \left[ \left(\frac{4}{5}\right)^t \right]_0^2 = 5 \left[ \left(\frac{4}{5}\right)^2 - \left(\frac{4}{5}\right)^0 \right] = -\frac{9}{5}.
    \]
    \item [β)] Ο όγκος της βενζίνης (σε λίτρα) που περιέχει το δοχείο δύο εβδομάδες μετά το άνοιγμα του καπακιού του δοχείου είναι ίσος με $g(2)$ ενώ αρχικά είναι $g(0)=5$.
    
    \OnPar{Α τρόπος}
    Είναι
    \[
        \int_0^2 5 \cdot \left(\frac{4}{5}\right)^t \cdot \ln\frac{4}{5} dt = -\frac{9}{5} \Leftrightarrow \int_0^2 g'(t)dt = -\frac{9}{5} \Leftrightarrow g(2)-g(0) = -\frac{9}{5} \Leftrightarrow g(2)-5 = -\frac{9}{5} \Leftrightarrow g(2) = \frac{16}{5}.
    \]
    
    \OnPar{Β τρόπος}
    Για κάθε $t>0$ έχουμε: $g'(t) = 5 \cdot \left(\frac{4}{5}\right)^t \cdot \ln\frac{4}{5} \Rightarrow g'(t) = \left( 5 \cdot \left(\frac{4}{5}\right)^t \right)'$.
    Οι συναρτήσεις $g(t)$ και $f(t) = 5 \cdot \left(\frac{4}{5}\right)^t$ είναι συνεχείς στο $[0,+\infty)$ και ισχύει $g'(t) = f'(t)$ για κάθε $t \in (0,+\infty)$, άρα υπάρχει σταθερά $c$ τέτοια, ώστε για κάθε $t \in [0,+\infty)$ να ισχύει
    \[
        g(t) = f(t)+c \Rightarrow g(t) = 5 \cdot \left(\frac{4}{5}\right)^t + c.
    \]
    Είναι $g(0)=5 \Rightarrow 5 \cdot \left(\frac{4}{5}\right)^0 + c = 5 \Rightarrow c=0$ οπότε $g(t) = 5 \cdot \left(\frac{4}{5}\right)^t$, $t \in [0,+\infty)$, άρα $g(2) = \frac{16}{5}$.

    \item [γ)] Είναι $\lim_{t\to+\infty} g(t) = \lim_{t\to+\infty} \left( 5 \cdot \left(\frac{4}{5}\right)^t \right) = 5 \lim_{t\to+\infty} \left(\frac{4}{5}\right)^t = 5 \cdot 0 = 0$ άρα καθώς ο χρόνος αυξάνεται απεριόριστα μόνο η μυρωδιά της βενζίνης θα υπάρχει στο δοχείο.
\end{alist}

\noindent \textbf{Σχόλιο:} Στην πράξη μετά από 40 εβδομάδες ο όγκος της βενζίνης του δοχείου είναι ίσος με
$g(40) = 5 \cdot \left(\frac{4}{5}\right)^{40} < 0,001$ που σημαίνει ότι το δοχείο δεν περιέχει ούτε καν 1 ml. Δηλαδή στην πράξη μετά από 40 εβδομάδες μπορούμε να πούμε ότι η βενζίνη έχει εξατμιστεί.